% Mục A.3 phần cài đặt. Phụ trách: TV2.
\phutrach{TV2 -- mục A.3, phần cài đặt}

Chương này chỉ trình bày môi trường và cách cài đặt ở mức tổng quát. Toàn bộ lệnh, file cấu hình, mã nguồn và kết quả đo nằm trong kho mã của nhóm (\url{https://github.com/taesikwoo268/nhom7-tidb-indexing}); muốn dựng lại môi trường, chỉ cần làm theo file \path{README.md}.

\subsection{Ba máy thực nghiệm}

Mỗi thành viên đo phần của mình trên máy cá nhân (Bảng~\ref{tab:ba-may}). Ba máy có phần cứng khác nhau, nên nhóm chỉ so số tuyệt đối giữa hai hệ trên cùng một máy; giữa các máy, nhóm chỉ so tỉ lệ.

\begin{table}[!htp]
\centering\bangnho
\caption{Ba máy thực nghiệm}\label{tab:ba-may}
\begin{tabular}{|C{1.1cm}|L{2.6cm}|L{3.2cm}|C{1.3cm}|L{2.4cm}|L{3.2cm}|}
\hline
\dongdau \textbf{Máy} & \textbf{Phần đo} & \textbf{CPU, số luồng} & \textbf{RAM} & \textbf{Ổ dữ liệu} & \textbf{Windows, Docker, WSL2} \\ \hline
M1 & RQ1 (3 triệu dòng, đối chiếu) & Intel Core i5-1235U, 10 nhân, 12 luồng & 16 GB & SSD NVMe 256 GB (KIOXIA) & Windows 11 Pro (10.0.26200); Docker 29.6.2; WSL2 10 GB RAM, 4 GB swap \\ \hline
M2 & RQ4, RQ3 (1 triệu dòng); RQ1, B.1 (1 và 3 triệu dòng) & Intel Core Ultra 5 125H, 14 nhân, 18 luồng & 16 GB & SSD NVMe 512 GB (UMIS) & Windows 11 Home (10.0.26300); Docker 29.8.0; WSL2 10 GB RAM, 4 GB swap \\ \hline
M3 & RQ1 (1 triệu dòng); RQ2 & Intel Core i5-13420H, 8 nhân, 12 luồng & 16 GB & SSD NVMe 512 GB (SK hynix) & Windows 11 Home Single Language (10.0.26200); Docker 29.8.2; WSL2 10 GB RAM, 4 GB swap \\ \hline
\end{tabular}
\end{table}

Với giới hạn WSL2 như trên, Docker thấy 12 CPU và 9,7 GiB bộ nhớ trên M1, 18 CPU và 9,7 GiB trên M2, 10 CPU và 9,7 GiB trên M3. Trong các lượt đo chính thức, máy luôn được cắm sạc để chế độ tiết kiệm pin không làm thay đổi hiệu năng.

\subsection{Cấu hình tài nguyên}

Cả ba máy dùng chung một hồ sơ tài nguyên mà nhóm gọi là XS (Bảng~\ref{tab:xs}). Hồ sơ này được thu nhỏ để cụm TiDB chạy được trên laptop 16 GB, trong khi hai hệ vẫn có cùng tổng giới hạn CPU và bộ nhớ đệm dữ liệu gần bằng nhau. Mỗi lúc chỉ có một hệ chạy.

\begin{table}[!htp]
\centering\bangfont
\caption{Hồ sơ XS: giới hạn tài nguyên của từng thành phần}\label{tab:xs}
\begin{tabular}{|L{4.2cm}|C{2.4cm}|C{3.0cm}|L{4.6cm}|}
\hline
\dongdau \textbf{Thành phần} & \textbf{CPU} & \textbf{RAM (MiB)} & \textbf{Bộ nhớ đệm} \\ \hline
PD & 0,4 & 512 & --- \\ \hline
TiDB server v8.5.5 & 1,2 & 1.280 & --- \\ \hline
TiKV v8.5.5 (3 node) & $3 \times 0{,}8$ & $3 \times 1.280$ & Block cache $3 \times 250$ MiB \\ \hline
\textbf{Tổng cụm TiDB} & \textbf{4,0} & \textbf{5.632} & \textbf{750 MiB} \\ \hline
MySQL 8.4.11 & 4,0 & 2.048 & Buffer pool 752 MiB (đặt 750M) \\ \hline
sysbench (RQ4) & 2,0 & 512 & --- \\ \hline
\end{tabular}
\end{table}

Hai hệ có cùng tổng 4 CPU và bộ nhớ đệm dữ liệu gần bằng nhau, nhưng tổng RAM khác nhau: cụm TiDB được 5,5 GiB, MySQL 2 GiB, vì cụm TiDB có năm tiến trình và tiến trình nào cũng cần bộ nhớ riêng. MySQL là một tiến trình nên dùng được cả 4 CPU, còn mỗi TiKV bị giới hạn ở 0,8 CPU kể cả khi hai TiKV kia rảnh; Chương~\ref{sec:ghi} cho thấy giới hạn này có ảnh hưởng thật. Buffer pool đặt 750M được MySQL làm tròn thành 752 MiB.

MySQL đặt \code{innodb\_flush\_log\_at\_trx\_commit=1} và tắt binlog; TiDB giữ cấu hình Raft log mặc định, mỗi TiKV có \code{memory-usage-limit} 1.000 MiB. Image TiKV chính thức được loại phần thông tin debug, mã thực thi giữ nguyên; nhờ vậy bộ nhớ một TiKV chiếm khi cụm rỗng giảm từ khoảng 1,9 GB xuống 0,16 GB. Về collation, MySQL dùng \code{utf8mb4\_unicode\_ci} còn TiDB dùng \code{utf8mb4\_bin}; chỉ \code{tracking\_code} là cột chuỗi. Mục~\ref{sec:rq4-collation} cho thấy khác biệt này không vượt mức nhiễu ở RQ4; RQ1 vẫn dùng collation mặc định nên ở Q1 chỉ so tỉ số trong từng hệ.

\subsection{Số node, số bản sao, quorum và khả năng chịu lỗi}

Số node TiKV và số bản sao của mỗi Region là hai khái niệm khác nhau, nên cần nói rõ cả hai. Trong đề tài, cả hai cấu hình TiDB đều có đúng 3 TiKV; nhóm chỉ đổi số bản sao qua tham số \code{replication.max-replicas} của PD (Bảng~\ref{tab:ban-sao}).

\begin{table}[!htp]
\centering\bangnho
\caption{Định nghĩa cấu hình bản sao}\label{tab:ban-sao}
\begin{tabular}{|C{1.4cm}|C{1.4cm}|C{2.0cm}|C{1.8cm}|C{1.7cm}|L{5.5cm}|}
\hline
\dongdau \textbf{Ký hiệu} & \textbf{TiKV} & \textbf{Bản sao mỗi Region} & \textbf{Quorum} & \textbf{Chịu mất} & \textbf{Dùng cho} \\ \hline
C3 & 3 & 3 & 2 trên 3 & 1 store & Mọi phép đo chính \\ \hline
C1 & 3 & 1 & 1 trên 1 & Không & Tách chi phí nhân bản ở RQ4, bộ 0 và 5 chỉ mục \\ \hline
\end{tabular}
\end{table}

Ở C3, một lệnh ghi được commit khi 2 trong 3 peer đã ghi Raft log~\cite{raft}, nên mất một store thì các Region vẫn đủ đa số để bầu leader mới; ở C1, store chứa peer duy nhất dừng thì Region cũng dừng. MySQL chạy một node, không có bản sao. Một TiKV duy nhất không thể tạo cấu hình 3 bản sao vì các peer của một Region phải nằm trên các store khác nhau. Ba TiKV ở đây là ba container trên một máy và một ổ đĩa, nên ``chịu mất một store'' chỉ có nghĩa là chịu mất một tiến trình; PD và TiDB server chỉ có một node nên vẫn là điểm lỗi đơn.

Trước mỗi lượt đo TiDB, chương trình kiểm tra \code{max-replicas}, đủ 3 store ở trạng thái \code{Up} và mọi Region có đúng số peer; thiếu một điều kiện là chương trình dừng.

\subsection{Các bước cài đặt}

Việc cài đặt gồm năm bước. Vì mọi cấu hình đều nằm trong Git, làm đúng các bước này thì ba máy sẽ có môi trường giống nhau.

\begin{enumerate}
  \item Cài Docker Desktop với WSL2, đặt \code{memory=10GB} và \code{swap=4GB} cho WSL2.
  \item Lấy mã nguồn về và tạo môi trường Python 3.10 trở lên.
  \item Khai báo tên máy và hồ sơ (\code{MACHINE\_ID}, \code{PROFILE=XS}). Nếu cổng 3306 đã bị chiếm thì đổi cổng MySQL; M2 dùng cổng 3307.
  \item Chạy lệnh kiểm tra. Lệnh này từ chối khi sai hồ sơ, Docker thiếu bộ nhớ hoặc CPU, máy đang chạy pin, hoặc có file cấu hình khác với Git.
  \item Bật từng hệ bằng \code{use\_engine.ps1} và nạp dữ liệu bằng \code{load.ps1}.
\end{enumerate}

\subsection{Minh chứng cài đặt}

Hình~\ref{fig:minh-chung-cai-dat} là ảnh chụp terminal trên M2 khi bật cụm TiDB: lệnh kiểm tra cho \code{team\_hash} chung của nhóm và kết quả \code{DAT}, năm container của cụm đang chạy, phiên bản v8.5.5 và ba store ở trạng thái Up với \code{max-replicas} bằng 3. Log nạp dữ liệu cho thấy hai hệ cùng nạp 1 triệu dòng với cùng tổng \code{cod\_amount} (5000968738439) và cùng mã SHA-256 của file nguồn; MySQL báo phiên bản 8.4.11, còn TiDB báo chuỗi phiên bản tương thích MySQL 8.0.11 kèm phiên bản thật v8.5.5.

\hinh{minh_chung_cai_dat.png}{0.8}{Ảnh chụp terminal trên M2: kiểm tra cấu hình, các container đang chạy, phiên bản TiDB và trạng thái ba store}{fig:minh-chung-cai-dat}
